\documentclass[11pt]{article}
\usepackage{../LaTeX/metricsstyle}
\usepackage{graphicx}

%% Assignment #1, ECON*6140 -- full solution, questions (1)-(5).
%% Numbers come from HW1.ipynb; the figure is ../LaTeX/hw01_q2_plot.pdf.
%% No bold anywhere: \question and the title block are redefined below.

\renewcommand{\question}[1]{\par\needspace{4\baselineskip}\medskip\noindent #1\par\nobreak\smallskip}

\begin{document}

\begin{center}
  {\footnotesize\scshape Department of Economics and Finance\\
   Lang School of Business and Economics}\\[5pt]
  {\large\scshape ECON*6140 Econometrics I}\\[3pt]
  {\LARGE Assignment \#1}\\[4pt]
  {\itshape Solution}
\end{center}
\vspace{0.4em}\hrule\vspace{0.9em}

\question{(1)}
Since $X = [X_{1}|X_{2}]$,
\[
  \begin{pmatrix} b_{1} \\ b_{2} \end{pmatrix}
  = (X\xT X)^{-1}X\xT y
  =
  \begin{pmatrix}
    X_{1}\xT X_{1} & X_{1}\xT X_{2} \\
    X_{2}\xT X_{1} & X_{2}\xT X_{2}
  \end{pmatrix}^{-1}
  \begin{pmatrix} X_{1}\xT y \\ X_{2}\xT y \end{pmatrix}.
\]
Let $M_{i} = I - X_{i}(X_{i}\xT X_{i})^{-1}X_{i}\xT$. Then
$X_{1}\xT X_{1} - X_{1}\xT X_{2}(X_{2}\xT X_{2})^{-1}X_{2}\xT X_{1} = X_{1}\xT M_{2}X_{1}$
and likewise for $X_{2}\xT M_{1}X_{2}$, so by the partitioned-inverse formula (3.71)
\[
  (X\xT X)^{-1}
  =
  \begin{pmatrix}
    (X_{1}\xT M_{2}X_{1})^{-1}
    & -(X_{1}\xT M_{2}X_{1})^{-1}X_{1}\xT X_{2}(X_{2}\xT X_{2})^{-1} \\[6pt]
    -(X_{2}\xT M_{1}X_{2})^{-1}X_{2}\xT X_{1}(X_{1}\xT X_{1})^{-1}
    & (X_{2}\xT M_{1}X_{2})^{-1}
  \end{pmatrix}.
\]
Multiplying out each row:
\begin{align*}
  b_{1} &= (X_{1}\xT M_{2}X_{1})^{-1}X_{1}\xT\bigl(I - X_{2}(X_{2}\xT X_{2})^{-1}X_{2}\xT\bigr)y
         = (X_{1}\xT M_{2}X_{1})^{-1}X_{1}\xT M_{2}\,y, \\
  b_{2} &= (X_{2}\xT M_{1}X_{2})^{-1}X_{2}\xT\bigl(I - X_{1}(X_{1}\xT X_{1})^{-1}X_{1}\xT\bigr)y
         = (X_{2}\xT M_{1}X_{2})^{-1}X_{2}\xT M_{1}\,y.
\end{align*}

\question{(2)}
Sample 1950:4--1996:4, $n = 185$. OLS gives (standard errors in brackets)
\[
  \widehat{\Delta r}_{t} = \underset{(0.1256)}{-0.2319}
  + \underset{(0.0200)}{0.0161}\,\pi_{t-1}
  + \underset{(5.7589)}{18.3807}\,\Delta y_{t-1}
  + \underset{(0.0741)}{0.2375}\,\Delta r_{t-1}
  - \underset{(0.0725)}{0.1540}\,\Delta r_{t-2},
  \qquad R^{2} = 0.134.
\]

\begin{center}
  \includegraphics[width=0.9\textwidth]{../LaTeX/hw01_q2_plot.pdf}\\
  {\small Fitted values and residuals, 1950:4--1996:4.}
\end{center}

Residuals on a constant and $\hat y$: both coefficients are zero ($10^{-16}$ is rounding) and $R^{2} = 0$.
This is an identity, not a finding. The normal equations give $X\xT e = 0$. The constant is a column of $X$
and $\hat y = Xb$, so $[\iota\ \hat y]\xT e = 0$, which is exactly the normal equation of this regression with solution $(0,0)$.

Fitted values on a constant and $e$: the slope is zero and $R^{2} = 0$, for the same reason,
since $\hat y\xT e = 0$ and $\bar e = 0$. The intercept is $\bar{\hat y} = \bar y = 0.0134$,
because $\sum e_{t} = 0$ makes the mean of $\hat y$ equal the mean of $\Delta r_{t}$.

\question{(3)}
Let $X_{1} = [\iota,\ \Delta y_{t-1},\ \Delta r_{t-1},\ \Delta r_{t-2}]$ and $x_{2} = \pi_{t-1}$.
Then $e_{t} = M_{1}\Delta r$, $e^{\pi}_{t} = M_{1}x_{2}$, and regressing $e_{t}$ on $e^{\pi}_{t}$ gives
\[
  \hat\beta = (x_{2}\xT M_{1}x_{2})^{-1}x_{2}\xT M_{1}\Delta r = 0.0161,
\]
the coefficient on $\pi_{t-1}$ in (1), as $b_{2}$ from question (1) says it must be.
The residuals are also the same as in (1): the largest difference is $9\times10^{-15}$.
This is the Frisch--Waugh--Lovell theorem.

The standard error is not quite the same, $0.0198$ against $0.0200$. The residual sum of squares is the same,
but the last regression divides it by $n - 1 = 184$ instead of $n - k = 180$.
Scaling by $\sqrt{184/180}$ recovers $0.0200$.

\question{(4)}
Sample 1953:1--1996:4, $n = 176$, with $c_{t} = \log CE_{t}$ and $y_{t} = \log YD_{t}$. OLS gives
\[
  \widehat{\Delta c}_{t} = \underset{(0.0012)}{0.0051}
  + \underset{(0.0545)}{0.3058}\,\Delta y_{t}
  + \underset{(0.0544)}{0.1242}\,\Delta y_{t-1}
  + \underset{(0.0551)}{0.0425}\,\Delta y_{t-2}
  + \underset{(0.0545)}{0.1239}\,\Delta y_{t-3}
  - \underset{(0.0539)}{0.1233}\,\Delta y_{t-4}.
\]

Method 1. Let $w = (0,1,1,1,1,1)\xT$. Then
\[
  \hat\gamma = w\xT\hat\beta = 0.4732, \qquad
  \operatorname{se}(\hat\gamma) = \sqrt{w\xT \widehat{\Var}(\hat\beta)\,w} = 0.1000.
\]
This is not the sum of the five standard errors. The covariances between the $\hat\beta_{i}$ enter through the off-diagonal terms.

Method 2. Substitute $\beta_{2} = \gamma - \beta_{3} - \beta_{4} - \beta_{5} - \beta_{6}$ into (2):
\[
  \Delta c_{t} = \beta_{1} + \gamma\,\Delta y_{t}
  + \beta_{3}(\Delta y_{t-1} - \Delta y_{t})
  + \beta_{4}(\Delta y_{t-2} - \Delta y_{t})
  + \beta_{5}(\Delta y_{t-3} - \Delta y_{t})
  + \beta_{6}(\Delta y_{t-4} - \Delta y_{t}) + \text{errors}.
\]
The coefficient on $\Delta y_{t}$ is now $\gamma$. OLS gives $\hat\gamma = 0.4732$ with standard error $0.1000$, the same as method 1.
The new regressors span the same space as the old ones, so the fit and the residuals do not change, only the parameters.

\question{(5)}
In the original regression $\hat\beta_{2} = (X_{2}\xT M_{1}X_{2})^{-1}X_{2}\xT M_{1}y$ and the residuals are $M_{X}y$.
Use $P_{1}X_{1} = X_{1}$, $M_{1}X_{1} = 0$, $P_{X}X_{i} = X_{i}$, and $P_{1}P_{X} = P_{1}$.

\begin{flatlist}
  \item[(a)] $\tilde\beta_{2} = (X_{2}\xT X_{2})^{-1}X_{2}\xT y$, residuals $M_{2}y$. Neither is the same, unless $X_{1}\xT X_{2} = 0$.
  \item[(b)] $\tilde\beta_{2} = (X_{2}\xT X_{2})^{-1}X_{2}\xT P_{1}y$, residuals $M_{2}P_{1}y$. Neither is the same.
  \item[(c)] $\tilde\beta_{2} = (X_{2}\xT P_{1}X_{2})^{-1}X_{2}\xT P_{1}y$, residuals $P_{1}y - P_{1}X_{2}\tilde\beta_{2}$. Neither is the same.
  \item[(d)] $\tilde\beta = (X\xT X)^{-1}X\xT P_{X}y = (X\xT X)^{-1}X\xT y$, so $\beta_{2}$ is the same.
             The residuals are $P_{X}y - P_{X}y = 0$, not the same.
  \item[(e)] $\tilde\beta_{2} = (X_{2}\xT X_{2})^{-1}X_{2}\xT P_{X}y = (X_{2}\xT X_{2})^{-1}X_{2}\xT y$, the same as (a). Neither is the same.
  \item[(f)] $\tilde\beta_{2} = (X_{2}\xT X_{2})^{-1}X_{2}\xT M_{1}y$, residuals $M_{2}M_{1}y$.
             The numerator is right but the denominator is $X_{2}\xT X_{2}$, not $X_{2}\xT M_{1}X_{2}$. Neither is the same.
  \item[(g)] $\tilde\beta_{2} = (X_{2}\xT M_{1}X_{2})^{-1}X_{2}\xT M_{1}y$, the same (FWL).
             Residuals $M_{1}y - M_{1}X_{2}\hat\beta_{2} = M_{1}(y - X_{1}\hat\beta_{1} - X_{2}\hat\beta_{2}) = M_{1}M_{X}y = M_{X}y$, the same.
  \item[(h)] $X_{1}\xT M_{1}X_{2} = 0$, so the two blocks are orthogonal and can be estimated separately.
             $\tilde\beta_{2}$ is as in (g), the same, and $\tilde\beta_{1} = (X_{1}\xT X_{1})^{-1}X_{1}\xT M_{1}y = 0$.
             Residuals the same as in (g), since $X_{1}$ contributes nothing to the fit.
  \item[(i)] $M_{1}X_{1} = 0$, so that block is a column of zeros and $\beta_{1}$ is not identified.
             Dropping it leaves (g): $\beta_{2}$ and the residuals are the same.
  \item[(j)] $M_{1}P_{X}y = M_{1}(X_{1}\hat\beta_{1} + X_{2}\hat\beta_{2}) = M_{1}X_{2}\hat\beta_{2}$, so
             $\tilde\beta_{2} = (X_{2}\xT M_{1}X_{2})^{-1}X_{2}\xT M_{1}X_{2}\hat\beta_{2} = \hat\beta_{2}$, the same.
             Residuals $P_{X}y - M_{1}X_{2}\hat\beta_{2} = P_{1}P_{X}y = P_{1}y$, not the same.
\end{flatlist}

So $\beta_{2}$ is the same in (d), (g), (h), (i) and (j), and the residuals are the same only in (g), (h) and (i).
In (g)--(i) the dependent variable is $M_{1}y$ and the fit removes $M_{1}X_{2}\hat\beta_{2}$ from it, which leaves $M_{X}y$.
I checked all ten numerically on the question (2) data with $X_{2} = \pi_{t-1}$.

\end{document}
